% ID: SEL-D0013-BA-E01
% ESTADO: validada. S014; tres pasadas del mismo autor, sin segundo revisor.
\textbf{a) Matrices con filas y columnas identificadas.}\par\nopagebreak[4]
Elegimos filas de $A$ para habitaciones individuales, dobles y triples, y columnas para agencias primera y segunda. En $D$, las filas son institutos primero, segundo y tercero, y las columnas los mismos tres tipos de habitación:
\[
A=\begin{pmatrix}65&78\\85&83\\104&106\end{pmatrix}_{3\times2},\qquad
D=\begin{pmatrix}3&15&2\\2&12&5\\1&16&7\end{pmatrix}_{3\times3}.
\]
Cada entrada de $A$ es un precio por habitación y noche; cada entrada de $D$ es una cantidad de habitaciones. Esta orientación permite el producto $DA$.

\textbf{b) Costes totales y elección.}\par\nopagebreak[4]
\[
DA=\begin{pmatrix}1678&1691\\1670&1682\\2153&2148\end{pmatrix}.
\]
Las filas representan institutos y las columnas agencias; las entradas son euros por noche para el total solicitado. Por ejemplo, la entrada primera es $3\cdot65+15\cdot85+2\cdot104=1678$.
Comparando cada fila, al primer instituto le interesa la \textbf{primera agencia}, con ahorro de $13$ euros/noche; al segundo, la \textbf{primera}, con ahorro de $12$; al tercero, la \textbf{segunda}, con ahorro de $5$.

\textbf{c) Invertibilidad.}\par\nopagebreak[4]
\[
\begin{aligned}
\det D&=3(12\cdot7-5\cdot16)-15(2\cdot7-5\cdot1)\\
&\qquad+2(2\cdot16-12\cdot1)\\
&=12-135+40=-83\ne0.
\end{aligned}
\]
Por tanto, \textbf{$D$ tiene inversa}. La matriz $A$ es rectangular $3\times2$, por lo que \textbf{no tiene inversa matricial ordinaria}. Otras orientaciones son posibles si se declaran y se adapta correctamente el orden del producto.
